% P15-U1: event limits, Borel--Cantelli and the Kolmogorov/Hewitt--Savage laws.
\chapter{Infinitely Often Events and Zero-One Laws}
\label{sec:p15-limsup-events-and-zero-one-laws}

An event may occur finitely often, eventually always, or infinitely often
without becoming permanent. These distinctions support the Borel--Cantelli
lemmas. Tail and finite-permutation invariance then give two different
reasons why a measurable event has probability zero or one.

The source is Durrett's author Version 5 dated 11 January 2019:
the event-limit definitions, Theorems 2.3.1 and 2.3.7, and the complete
pairwise counting result 2.3.9 \cite[pp.~67--72]{DurrettPTEAuthor2019};
and the tail/permutation definitions, examples and Theorems 2.5.3--2.5.4
\cite[pp.~81--83]{DurrettPTEAuthor2019}. The finite-coordinate
approximation needed for Hewitt--Savage is proved here directly.
No conditional-expectation or martingale theorem is used.

\section{Limits of events and Borel--Cantelli}
\label{sec:p15-event-limits}

\begin{hbdefinition}{def:p15-event-limsup-liminf}{Infinitely often and eventually}
On a probability space \((\Omega,\mathcal F,\mathbb P)\), let
\(A_1,A_2,\ldots\) be events. Define
\[
 \limsup_n A_n=\bigcap_{m\geq1}\bigcup_{n\geq m}A_n,
 \qquad
 \liminf_n A_n=\bigcup_{m\geq1}\bigcap_{n\geq m}A_n.
\]
The first is the event that \(A_n\) occurs \emph{infinitely often},
abbreviated \(A_n\text{ i.o.}\). The second is the event that
\(A_n\) occurs \emph{eventually}, meaning at every sufficiently
large index. These are pointwise statements about an outcome's memberships.
Let \(N_n=\sum_{j=1}^n\mathbf1_{A_j}\) count occurrences through
time \(n\), and let \(N_\infty=\sum_{j\geq1}\mathbf1_{A_j}\)
be the total count, possibly infinite. The count \(N_\infty\) is
infinite exactly on \(\limsup_n A_n\).
\index{limsup!events}\index{liminf!events}\index{infinitely often}
\end{hbdefinition}

\begin{hbtheorem}{thm:p15-event-limit-calculus}{proved=annex:p15.event-limit-calculus}{Event limits, indicators and almost-sure convergence}
The two event limits are measurable, and
\[
 \begin{aligned}
 (\limsup_n A_n)^c&=\liminf_n A_n^c,&
 (\liminf_n A_n)^c&=\limsup_n A_n^c,\\
 \mathbf1_{\limsup_n A_n}&=\limsup_n\mathbf1_{A_n},&
 \mathbf1_{\liminf_n A_n}&=\liminf_n\mathbf1_{A_n}.
 \end{aligned}
\]
Moreover,
\[
 \mathbb P(\limsup_n A_n)\geq\limsup_n\mathbb P(A_n),\qquad
 \mathbb P(\liminf_n A_n)\leq\liminf_n\mathbb P(A_n).
\]
For real random variables \(X_n,X\), convergence \(X_n\to X\)
a.s. is equivalent to
\[
 \mathbb P(|X_n-X|>1/r\text{ i.o.})=0
                        \qquad\text{for every integer }r\geq1.
\]
It is also equivalent to the same condition for every real threshold
\(\varepsilon>0\); the countable formulation supplies one common
probability-one event on which convergence holds.
\end{hbtheorem}

\begin{hbtheorem}{thm:p15-borel-cantelli-first}{proved=annex:p15.borel-cantelli-first}{The first Borel--Cantelli lemma}
For any event sequence, with no independence assumption,
\[
 \sum_{n\geq1}\mathbb P(A_n)<\infty
               \quad\Longrightarrow\quad
 \mathbb P(A_n\text{ i.o.})=0.
\]
More precisely,
\(\mathbb E N_\infty=\sum_n\mathbb P(A_n)\), and
\(\mathbb P(\bigcup_{n\geq m}A_n)\leq\sum_{n\geq m}\mathbb P(A_n)\).
In particular, if
\(\sum_n\mathbb P(|X_n-X|>1/r)<\infty\) for each positive integer
\(r\), then \(X_n\to X\) a.s.
\end{hbtheorem}

\begin{hbtheorem}{thm:p15-borel-cantelli-second}{proved=annex:p15.borel-cantelli-second}{The source's second Borel--Cantelli lemma}
If \((A_n)\) are mutually independent events and
\(\sum_n\mathbb P(A_n)=\infty\), then
\(\mathbb P(A_n\text{ i.o.})=1\).
For fixed \(m\leq N\), the proof gives the explicit bound
\[
 \mathbb P\left(\bigcap_{n=m}^N A_n^c\right)
   =\prod_{n=m}^N(1-\mathbb P(A_n))
   \leq\exp\left(-\sum_{n=m}^N\mathbb P(A_n)\right).
\]
Thus for mutually independent events the finite/infinite probability
sum determines whether infinitely many occurrences have probability
zero or one. The next theorem establishes the separate stronger
pairwise-independent counting conclusion in the source.
\end{hbtheorem}

\begin{hbtheorem}{thm:p15-pairwise-event-count-law}{proved=annex:p15.pairwise-event-count-law}{The pairwise-independent occurrence-count law}
Let \((A_n)\) be pairwise independent and put
\(t_n=\sum_{j=1}^n\mathbb P(A_j)\). If \(t_n\to\infty\), then
\[
                       \frac{N_n}{t_n}\longrightarrow1\quad\text{a.s.}
\]
The quotient is used for all sufficiently large \(n\), when \(t_n>0\).
In particular \(\mathbb P(A_n\text{ i.o.})=1\) under pairwise
independence as well. The proof of Durrett's Theorem 2.3.9 uses
a variance bound, the first Borel--Cantelli lemma on a selected
subsequence, and interpolation between successive selected indices.
\end{hbtheorem}

\section{Finite information and the tail sigma-algebra}
\label{sec:p15-tail-information}

\begin{hbdefinition}{def:p15-tail-sigma-algebra}{Finite observations, the remote future and symmetric difference}
For a sequence of random elements \(X_1,X_2,\ldots\), define
\[
 \begin{aligned}
 \mathcal H_n&=\sigma(X_1,\ldots,X_n),&
 \mathcal H_\infty&=\sigma(X_1,X_2,\ldots),\\
 \mathcal F'_n&=\sigma(X_n,X_{n+1},\ldots),&
 \mathcal T&=\bigcap_{n\geq1}\mathcal F'_n.
 \end{aligned}
\]
Each sigma-algebra generated by several observations means the
sigma-algebra generated by the union of their individual generated
sigma-algebras. The field \(\mathcal T\) is the \emph{tail sigma-algebra}.
It records membership in every remote-future field. It is contained
in \(\mathcal H_\infty\), which need not equal the entire ambient
\(\mathcal F\).
For events \(A,B\), their \emph{symmetric difference} is
\(A\triangle B=(A\setminus B)\cup(B\setminus A)\).
\index{tail sigma-algebra}\index{symmetric difference}
\end{hbdefinition}

\begin{hbtheorem}{thm:p15-finite-coordinate-approximation}{proved=annex:p15.finite-coordinate-approximation}{Approximation by finite information}
Let \((\mathcal H_n)\) be increasing sub-sigma-algebras of a
probability space and let \(\mathcal H_\infty=\sigma(\bigcup_n\mathcal H_n)\).
For every \(A\in\mathcal H_\infty\) and \(\varepsilon>0\),
there exist \(n\) and \(B\in\mathcal H_n\) with
\(\mathbb P(A\triangle B)<\varepsilon\).
One can choose a full sequence \(A_n\in\mathcal H_n\) such that
\(\mathbb P(A_n\triangle A)\to0\).
The estimates used below are
\[
 |\mathbb P(A)-\mathbb P(B)|\leq\mathbb P(A\triangle B),
 \qquad
 A\triangle C\subseteq(A\triangle B)\cup(B\triangle C).
\]
On a canonical coordinate product \(S^{\mathbb N_+}\), an event
in \(\sigma(X_1,\ldots,X_n)\) has the form
\(\{\omega:(\omega_1,\ldots,\omega_n)\in B_n\}\),
where \(B_n\) belongs to the finite product value sigma-algebra.
\end{hbtheorem}

\begin{hbtheorem}{thm:p15-tail-event-structure}{proved=annex:p15.tail-event-structure}{Tail events and the precise finite-change criterion}
The field \(\mathcal T\) is a sigma-algebra. If \(B_n\) are
measurable value sets, then \(\{X_n\in B_n\text{ i.o.}\}\in\mathcal T\).
Taking \(X_n=\mathbf1_{A_n}\) and \(B_n=\{1\}\) gives
\(A_n\text{ i.o.}\) as a tail event of the indicator sequence.
For real \(X_n\), put \(S_n=\sum_{i=1}^nX_i\). The event that
\(S_n\) converges to a finite real value is a tail event, as are
the events of convergence to \(+\infty\) or to \(-\infty\).
If deterministic \(c_n\to+\infty\), then for every real \(a\),
\[
 \{\limsup_n S_n/c_n>a\},\qquad
 \{\limsup_n S_n/c_n\geq a\}
                         \quad\text{belong to }\mathcal T.
\]
Ratios are taken after the index at which all \(c_n>0\);
finitely many initial terms do not affect the limit superior.

On the canonical measurable product \(S^{\mathbb N_+}\) of a
nonempty measurable space, a product-measurable event is in the tail
sigma-algebra exactly when membership is unchanged by every change
to finitely many coordinates. This pointwise criterion refers to
the canonical product, not to arbitrary modifications of an unspecified
underlying sample space.
On \(\mathbb R^{\mathbb N_+}\), the event
\(\{\limsup_n S_n>0\}\) is not a tail event.
\end{hbtheorem}

\begin{hbtheorem}{thm:p15-kolmogorov-tail-zero-one}{proved=annex:p15.kolmogorov-tail-zero-one}{Kolmogorov's tail zero-one law}
If the sequence \((X_n)\) is mutually independent, then
\(\mathcal T\) is independent of \(\mathcal H_\infty\).
In particular every \(A\in\mathcal T\) has
\(\mathbb P(A)\in\{0,1\}\).
Identical distributions and moment assumptions are unnecessary.
Pairwise independence alone is not a hypothesis of this theorem.
\index{Kolmogorov zero-one law}
\end{hbtheorem}

\section{Finite permutations and the Hewitt--Savage law}
\label{sec:p15-hewitt-savage}

\begin{hbdefinition}{def:p15-permutable-events}{Permutable events and the exchangeable sigma-algebra}
Let \(S\) be a nonempty measurable space with value sigma-algebra
\(\mathcal S\). Work on \(\Omega=S^{\mathbb N_+}\) with its
coordinate product sigma-algebra and coordinate maps \(X_i(\omega)=\omega_i\).
A \emph{finite permutation} \(\pi\) is a bijection of
\(\mathbb N_+\) equal to the identity outside a finite set. Define
\(T_\pi(\omega)_i=\omega_{\pi(i)}\).
An event is \emph{permutable} when
\(T_\pi^{-1}(A)=A\) for every finite permutation. Let
\(\mathcal E\) be the collection of product-measurable permutable
events, called the \emph{exchangeable sigma-algebra}.

This is exact equality of events. A probability law is itself
\emph{exchangeable} when \(\mathbb P(T_\pi^{-1}(B))=\mathbb P(B)\)
for every product-measurable \(B\) and every finite permutation.
Invariance of an event and invariance of a probability law are distinct
conditions. Neither definition assumes independence.
\index{permutable event}\index{exchangeable sigma-algebra}\index{exchangeable law}
\end{hbdefinition}

\begin{hbtheorem}{thm:p15-permutation-structure}{proved=annex:p15.permutation-structure}{Measurable permutations, tail inclusion and source examples}
Every \(T_\pi\) is a measurable bijection with measurable inverse,
\(\mathcal E\) is a sigma-algebra, and \(\mathcal T\subseteq\mathcal E\).
If the coordinate sequence is i.i.d. under a given probability law,
that law is invariant under every finite permutation.

For \(S=\mathbb R\), finite permutations leave \(S_n\)
unchanged at all sufficiently large \(n\). Consequently
\(\{S_n\in B\text{ i.o.}\}\) is permutable for every Borel \(B\).
So are threshold events for \(\limsup_n S_n/c_n\) whenever the
deterministic \(c_n\) are nonzero for all sufficiently large \(n\).
The inclusion of the tail sigma-algebra need not be equality as
collections of events, even under an i.i.d. probability law.
\end{hbtheorem}

\begin{hbtheorem}{thm:p15-hewitt-savage-zero-one}{proved=annex:p15.hewitt-savage-zero-one}{The Hewitt--Savage zero-one law}
On the canonical product above, suppose the coordinate sequence is
independent and identically distributed. Then every
\(A\in\mathcal E\) has \(\mathbb P(A)\in\{0,1\}\).
For an i.i.d. sequence on another probability space, the pullback of
any such product-measurable permutable event has probability zero or one.

Under an i.i.d. law, both \(\mathcal T\) and \(\mathcal E\)
are therefore \emph{probability-trivial}: each of their events has
probability zero or one. This does not identify their literal sets
of events. The proof supplies finite-coordinate approximations and
the source's independent-block swap in full.
\index{Hewitt--Savage zero-one law}
\end{hbtheorem}

\begin{hbcorollary}{thm:p15-zero-one-measurable-values}{proved=annex:p15.zero-one-measurable-values}{Measurable values on a probability-trivial sigma-algebra}
If every event in a sub-sigma-algebra \(\mathcal G\) has probability
zero or one, every \(\mathcal G\)-measurable extended real variable
\(Z\) equals one constant \(c\in[-\infty,\infty]\) a.s.
If \(Z\) is finite a.s., then \(c\) is real. This applies to
tail-measurable variables under mutual independence and to
\(\mathcal E\)-measurable variables under an i.i.d. law.
It does not determine the constant's value or by itself prove
convergence of a sequence.
\end{hbcorollary}

\section{Worked examples and failed hypotheses}
\label{sec:p15-zero-one-boundaries}

\begin{hbtheorem}{thm:p15-zero-one-countermodels}{proved=annex:p15.zero-one-countermodels}{Rare events, dependence and distinct trivial sigma-algebras}
\par\nobreak
\begin{enumerate}
\item Independent binary coordinates with \(\mathbb P(X_n=1)=1/n\)
      exist by the countable product construction. They equal one
      infinitely often a.s. If \(N_n\) counts these occurrences
      and \(H_n=\sum_{j=1}^n1/j\), then \(N_n/H_n\to1\) a.s.
      In contrast, probabilities \(1/n^2\) give only finitely many
      ones a.s.
\item On the uniform space \((0,1)\), let \(0\leq a_n\leq1\),
      \(a_n\to0\), and \(A_n=(0,a_n)\). Then
      \(\limsup_n A_n=\varnothing\), although the probability
      sum diverges whenever \(a_n\geq1/n\). In particular
      \(a_n=1/n\) gives the stated failure. This is the source's
      shrinking-interval counterexample.
\item Let \(Y,B_1,B_2,\ldots\) be mutually independent fair bits,
      and define \(X_{2j-1}=B_j\), \(X_{2j}=Y\oplus B_j\),
      where \(u\oplus v\) is addition modulo two on \(\{0,1\}\).
      The \(X_i\) are identically distributed and pairwise independent,
      but \(\{Y=1\}\) belongs to their tail sigma-algebra and
      has probability \(1/2\). Pairwise independence does not suffice
      for Kolmogorov's zero-one law.
\item On the canonical product of \(\{0,1\}\) with all subsets
      measurable, the event
      \(A=\{\text{at least one coordinate equals one}\}\) is
      permutable but not tail. Under i.i.d. fair bits it has probability
      one, proving strict inclusion of the two sigma-algebras as sets.
      Under independent coordinates with only \(X_1\) a fair bit
      and every later coordinate zero a.s., it has probability \(1/2\).
      Thus identical distributions cannot be dropped from Hewitt--Savage.
\item If every coordinate equals one shared fair bit \(Y\), their
      joint law is invariant under all finite permutations and all
      coordinate laws agree. Nevertheless the tail and permutable event
      that ones occur infinitely often has probability \(1/2\).
      Exchangeability of the law without independence is insufficient.
\end{enumerate}
\end{hbtheorem}

\begin{hbmethod}{def:p15-zero-one-method}{Choosing the applicable infinite-event argument}
For an infinitely-often event, first examine the series of event
probabilities. A finite series invokes the first Borel--Cantelli lemma
without independence. A divergent series with mutual independence invokes
the source's second proof; pairwise independence is sufficient through
the separate occurrence-count theorem. For a tail event under mutually
independent observations, Kolmogorov gives a zero-or-one conclusion
without identical laws. For an event exactly invariant under finite
permutations and an i.i.d. law, Hewitt--Savage also applies beyond the tail
sigma-algebra. These conclusions concern probability, and do not assert
that every individual sample path has the asserted behavior.
\end{hbmethod}

\begin{proofboundary}{def:p15-zero-one-next}{Further limit and random-series arguments}
The two Borel--Cantelli lemmas, the pairwise counting refinement and both
zero-one laws are now supplied with their actual approximation prerequisites.
The source's subsequence/continuous-mapping applications, its fourth-moment
strong law, infinite-mean obstruction, record-value and longest-run
applications require their separate complete treatments. The next
random-series component starts at Kolmogorov's maximal inequality,
Theorem 2.5.5, then variance-summable convergence and the three-series
theorem. No such convergence theorem is inferred here merely from an
event's zero-or-one probability.
\end{proofboundary}
